% ============================================================================
%  preamble/lang-thai.tex  --  Thai language support
% ----------------------------------------------------------------------------
%  Load this in any deck that contains Thai. Load it AFTER preamble/base.tex.
%
%  ---------------------------------------------------------------------------
%  THE RULE: just type Thai. There is nothing to wrap.
%  ---------------------------------------------------------------------------
%      \title{มาตรฐานข้อมูลเปิดระดับ 5 ดาว}          works
%      \begin{frame}{ผลลัพธ์}                          works
%      ข้อความภาษาไทยธรรมดา                            works
%
%  On XeLaTeX the ucharclasses package below watches the character stream and
%  switches to Sarabun the moment a Thai codepoint appears, then switches back
%  to Fira Sans afterwards. You never call a macro.
%
%  ---------------------------------------------------------------------------
%  WHY THIS FILE EXISTS AT ALL
%  ---------------------------------------------------------------------------
%  Fira Sans, the deck's Latin typeface, contains no Thai glyphs. XeTeX does
%  NOT fall back to another font on its own -- it silently drops the character
%  and prints nothing, with no error and (by default) no warning. A slide can
%  come out blank while the log says the build succeeded.
%
%  That is the failure this file prevents.
%
%  ---------------------------------------------------------------------------
%  WHAT EACH ENGINE GIVES YOU
%  ---------------------------------------------------------------------------
%    XeLaTeX    real Sarabun, automatic, no wrapping        <- use this one
%    pdfLaTeX   Garuda for the whole document; Sarabun is a .ttf and pdfLaTeX
%               cannot read .ttf at all. A warning is printed so the
%               substitution is never silent.
%    LuaLaTeX   NOT SUPPORTED for Thai -- stops with an error telling you to
%               use XeLaTeX. See the LuaLaTeX block near the bottom.
% ============================================================================

% Safety net: main.tex normally works this out. Defined here too so the file
% still behaves if it is \input before the bootstrap for some reason.
\providecommand{\deckroot}{}

\ifPDFTeX
  % =========================================================================
  %  pdfLaTeX branch -- Garuda fallback
  % =========================================================================
  \GenericWarning{}{%
    ^^J**********************************************************^^J%
    * lang-thai: building with pdfLaTeX.                      *^^J%
    * Thai will render in GARUDA, not Sarabun.                *^^J%
    * Sarabun is a .ttf font and pdfLaTeX cannot read .ttf.   *^^J%
    * Build with XeLaTeX to get real Sarabun.                 *^^J%
    **********************************************************^^J%
  }

  % ------------------------------------------------------------------------
  %  NOTE: babel's `thai` option is deliberately NOT used here.
  %
  %  babel-thai predates UTF-8. Loading it sets bytes 0xA1-0xFB to catcode 11
  %  so that TIS-620 (an 8-bit Thai encoding from the 1990s) can be typed
  %  directly. That switches OFF LaTeX's UTF-8 decoder, and a modern UTF-8
  %  source file then fails with "Invalid UTF-8 byte" on every Thai word.
  %
  %  So we take only the two pieces we actually need -- the LTH font encoding
  %  and a Thai font -- and supply the UTF-8 mapping ourselves.
  % ------------------------------------------------------------------------
  \usepackage[LTH,T1]{fontenc}
  \usepackage[garuda]{fonts-tlwg}

  % Maps Thai Unicode codepoints onto their LTH font slots.
  % ============================================================================
%  thai-utf8-pdftex.def  --  read UTF-8 Thai under pdfLaTeX
% ----------------------------------------------------------------------------
%  GENERATED FILE -- mechanical, but safe to read.
%
%  The problem this solves:
%    The Thai support in TeX Live (babel-thai / fonts-tlwg) is built around
%    TIS-620, an 8-bit Thai encoding from the 1990s, and ships no UTF-8
%    mapping. pdfLaTeX reading a modern UTF-8 .tex file therefore stops with
%    "Invalid UTF-8 byte" on every Thai word.
%
%  The fix:
%    TIS-620 and Unicode agree on the order of Thai characters and differ by
%    a constant offset:
%
%        TIS-620 byte  =  Unicode codepoint  -  0x0E00  +  0xA0
%
%    So each Thai codepoint maps to its LTH font slot with one
%    \DeclareUnicodeCharacter line. That is what follows.
%
%  Two details that matter:
%
%    \thaislot is declared ROBUST. Thai characters often appear in frame
%    titles and section names, which LaTeX copies into the .toc/.nav files.
%    A fragile command breaks there with "Argument of \@secondoffive has an
%    extra }".
%
%    It expands to a bare \char, not to a wrapped box or a per-character font
%    switch. \char still takes part in TeX ligature and kerning, which is how
%    the LTH fonts stack Thai upper vowels and tone marks. Wrapping each
%    character individually would break that positioning.
%
%  Only loaded by preamble/lang-thai.tex, and only under pdfLaTeX.
% ============================================================================

\DeclareRobustCommand{\thaislot}[1]{\char#1\relax}

% --- Consonants, vowels and tone marks: U+0E01-U+0E3A ---------------------
\DeclareUnicodeCharacter{0E01}{\thaislot{161}}
\DeclareUnicodeCharacter{0E02}{\thaislot{162}}
\DeclareUnicodeCharacter{0E03}{\thaislot{163}}
\DeclareUnicodeCharacter{0E04}{\thaislot{164}}
\DeclareUnicodeCharacter{0E05}{\thaislot{165}}
\DeclareUnicodeCharacter{0E06}{\thaislot{166}}
\DeclareUnicodeCharacter{0E07}{\thaislot{167}}
\DeclareUnicodeCharacter{0E08}{\thaislot{168}}
\DeclareUnicodeCharacter{0E09}{\thaislot{169}}
\DeclareUnicodeCharacter{0E0A}{\thaislot{170}}
\DeclareUnicodeCharacter{0E0B}{\thaislot{171}}
\DeclareUnicodeCharacter{0E0C}{\thaislot{172}}
\DeclareUnicodeCharacter{0E0D}{\thaislot{173}}
\DeclareUnicodeCharacter{0E0E}{\thaislot{174}}
\DeclareUnicodeCharacter{0E0F}{\thaislot{175}}
\DeclareUnicodeCharacter{0E10}{\thaislot{176}}
\DeclareUnicodeCharacter{0E11}{\thaislot{177}}
\DeclareUnicodeCharacter{0E12}{\thaislot{178}}
\DeclareUnicodeCharacter{0E13}{\thaislot{179}}
\DeclareUnicodeCharacter{0E14}{\thaislot{180}}
\DeclareUnicodeCharacter{0E15}{\thaislot{181}}
\DeclareUnicodeCharacter{0E16}{\thaislot{182}}
\DeclareUnicodeCharacter{0E17}{\thaislot{183}}
\DeclareUnicodeCharacter{0E18}{\thaislot{184}}
\DeclareUnicodeCharacter{0E19}{\thaislot{185}}
\DeclareUnicodeCharacter{0E1A}{\thaislot{186}}
\DeclareUnicodeCharacter{0E1B}{\thaislot{187}}
\DeclareUnicodeCharacter{0E1C}{\thaislot{188}}
\DeclareUnicodeCharacter{0E1D}{\thaislot{189}}
\DeclareUnicodeCharacter{0E1E}{\thaislot{190}}
\DeclareUnicodeCharacter{0E1F}{\thaislot{191}}
\DeclareUnicodeCharacter{0E20}{\thaislot{192}}
\DeclareUnicodeCharacter{0E21}{\thaislot{193}}
\DeclareUnicodeCharacter{0E22}{\thaislot{194}}
\DeclareUnicodeCharacter{0E23}{\thaislot{195}}
\DeclareUnicodeCharacter{0E24}{\thaislot{196}}
\DeclareUnicodeCharacter{0E25}{\thaislot{197}}
\DeclareUnicodeCharacter{0E26}{\thaislot{198}}
\DeclareUnicodeCharacter{0E27}{\thaislot{199}}
\DeclareUnicodeCharacter{0E28}{\thaislot{200}}
\DeclareUnicodeCharacter{0E29}{\thaislot{201}}
\DeclareUnicodeCharacter{0E2A}{\thaislot{202}}
\DeclareUnicodeCharacter{0E2B}{\thaislot{203}}
\DeclareUnicodeCharacter{0E2C}{\thaislot{204}}
\DeclareUnicodeCharacter{0E2D}{\thaislot{205}}
\DeclareUnicodeCharacter{0E2E}{\thaislot{206}}
\DeclareUnicodeCharacter{0E2F}{\thaislot{207}}
\DeclareUnicodeCharacter{0E30}{\thaislot{208}}
\DeclareUnicodeCharacter{0E31}{\thaislot{209}}
\DeclareUnicodeCharacter{0E32}{\thaislot{210}}
\DeclareUnicodeCharacter{0E33}{\thaislot{211}}
\DeclareUnicodeCharacter{0E34}{\thaislot{212}}
\DeclareUnicodeCharacter{0E35}{\thaislot{213}}
\DeclareUnicodeCharacter{0E36}{\thaislot{214}}
\DeclareUnicodeCharacter{0E37}{\thaislot{215}}
\DeclareUnicodeCharacter{0E38}{\thaislot{216}}
\DeclareUnicodeCharacter{0E39}{\thaislot{217}}
\DeclareUnicodeCharacter{0E3A}{\thaislot{218}}

% --- Baht sign, remaining vowels, digits: U+0E3F-U+0E5B -------------------
%     (U+0E3B-U+0E3E are unassigned in Unicode, and have no TIS-620 byte)
\DeclareUnicodeCharacter{0E3F}{\thaislot{223}}
\DeclareUnicodeCharacter{0E40}{\thaislot{224}}
\DeclareUnicodeCharacter{0E41}{\thaislot{225}}
\DeclareUnicodeCharacter{0E42}{\thaislot{226}}
\DeclareUnicodeCharacter{0E43}{\thaislot{227}}
\DeclareUnicodeCharacter{0E44}{\thaislot{228}}
\DeclareUnicodeCharacter{0E45}{\thaislot{229}}
\DeclareUnicodeCharacter{0E46}{\thaislot{230}}
\DeclareUnicodeCharacter{0E47}{\thaislot{231}}
\DeclareUnicodeCharacter{0E48}{\thaislot{232}}
\DeclareUnicodeCharacter{0E49}{\thaislot{233}}
\DeclareUnicodeCharacter{0E4A}{\thaislot{234}}
\DeclareUnicodeCharacter{0E4B}{\thaislot{235}}
\DeclareUnicodeCharacter{0E4C}{\thaislot{236}}
\DeclareUnicodeCharacter{0E4D}{\thaislot{237}}
\DeclareUnicodeCharacter{0E4E}{\thaislot{238}}
\DeclareUnicodeCharacter{0E4F}{\thaislot{239}}
\DeclareUnicodeCharacter{0E50}{\thaislot{240}}
\DeclareUnicodeCharacter{0E51}{\thaislot{241}}
\DeclareUnicodeCharacter{0E52}{\thaislot{242}}
\DeclareUnicodeCharacter{0E53}{\thaislot{243}}
\DeclareUnicodeCharacter{0E54}{\thaislot{244}}
\DeclareUnicodeCharacter{0E55}{\thaislot{245}}
\DeclareUnicodeCharacter{0E56}{\thaislot{246}}
\DeclareUnicodeCharacter{0E57}{\thaislot{247}}
\DeclareUnicodeCharacter{0E58}{\thaislot{248}}
\DeclareUnicodeCharacter{0E59}{\thaislot{249}}
\DeclareUnicodeCharacter{0E5A}{\thaislot{250}}
\DeclareUnicodeCharacter{0E5B}{\thaislot{251}}


  % ------------------------------------------------------------------------
  %  WHY THE WHOLE DOCUMENT SWITCHES TO GARUDA
  %
  %  pdfTeX has no way to change font part-way through a run of characters --
  %  that trick (ucharclasses, in the XeLaTeX branch below) needs XeTeX. And
  %  Garuda ships ONLY in the LTH encoding, so it cannot be slotted in beside
  %  Fira Sans as a second family.
  %
  %  So the choice is: put the entire document in Garuda, or leave Thai to be
  %  typeset with Fira Sans's LTH slots -- which produces Latin gibberish
  %  ("ćéĺďÇÒÁÀÒÉÒäůÂ") rather than Thai. Garuda covers ASCII Latin perfectly
  %  well, so the whole document goes Garuda and stays readable.
  %
  %  This is the degraded path. Build with XeLaTeX and you keep Fira Sans for
  %  Latin and get real Sarabun for Thai.
  % ------------------------------------------------------------------------
  \AtBeginDocument{\thaitext}

  % Thai has no spaces between words, so give TeX slack rather than let it
  % push text into the margin.
  \sloppy

  % --- Compatibility aliases ----------------------------------------------
  %  Nothing needs wrapping any more. These stay so that older slides written
  %  against the previous version of this template keep working unchanged.
  \newenvironment{thaipar}{\par}{\par}
  \newcommand{\thaiinline}[1]{#1}
  \providecommand{\textthai}[1]{#1}

\else
  % =========================================================================
  %  XeLaTeX / LuaLaTeX branch -- real Sarabun
  % =========================================================================
  %  Sarabun is loaded BY FILE PATH, not by system font name, so the font does
  %  not need to be installed. It travels with the repo: a fresh copy of this
  %  folder, or an upload to Overleaf, both work.
  %
  %  \deckroot is worked out by main.tex and points at the folder holding
  %  theme/. Everything stays relative, so the deck can be moved anywhere.
  %
  %  Scale=MatchLowercase: Sarabun runs large for its point size next to
  %  Latin. This matches the two so Thai does not tower over English.
  % -------------------------------------------------------------------------
  \newfontfamily\thaifont[
    Path           = \deckroot theme/assets/fonts/,
    Extension      = .ttf,
    UprightFont    = Sarabun-Regular,
    ItalicFont     = Sarabun-Italic,
    BoldFont       = Sarabun-Bold,
    BoldItalicFont = Sarabun-BoldItalic,
    Script         = Thai,
    Scale          = MatchLowercase,
  ]{Sarabun}

  \let\thaifontsf\thaifont

  % -------------------------------------------------------------------------
  %  LINE SPACING
  %  Thai stacks up to two marks above the baseline (sara + tone, e.g. กื๊อ)
  %  and hangs vowels below it (ปรุง). Latin line spacing is not tall enough
  %  for that: the tone marks of one line collide with the descenders of the
  %  line above. This opens the lines up just enough to clear both.
  % -------------------------------------------------------------------------
  \linespread{1.25}

  \makeatletter
  %  Remember the family name fontspec generated for Sarabun, so the switch
  %  below can change ONLY the family and leave weight and shape alone. That
  %  is what keeps \textbf{ผลลัพธ์} bold instead of resetting it to regular.
  \newcommand*{\deckthaifam}{}
  {\thaifont\xdef\deckthaifam{\f@family}}
  \newcommand*{\decklatinfam}{}

  \newcommand{\deckenterthai}{%
    \xdef\decklatinfam{\f@family}%
    \fontfamily{\deckthaifam}\selectfont
    \ifxetex
      %  Thai is written without spaces between words. Left alone, TeX treats
      %  a whole clause as one unbreakable word, overflows the slide, and then
      %  inserts a HYPHEN mid-word -- which Thai never uses. Handing the line
      %  breaking to ICU's Thai dictionary fixes both.
      \XeTeXlinebreaklocale "th"%
      \XeTeXlinebreakskip=0pt plus 0.1em\relax
    \fi
    %  Belt and braces: no Latin hyphenation patterns inside Thai.
    \@ifundefined{l@nohyphenation}{}{\language\l@nohyphenation}%
  }

  \newcommand{\deckleavethai}{%
    \fontfamily{\decklatinfam}\selectfont
    \ifxetex
      \XeTeXlinebreaklocale "en"%
      \XeTeXlinebreakskip=0pt\relax
    \fi
    \language\l@english
  }
  \makeatother

  \ifxetex
    % -----------------------------------------------------------------------
    %  THE AUTOMATIC BIT (XeLaTeX only)
    %
    %  ucharclasses assigns every character to its Unicode block and lets us
    %  run code when the text crosses from one block into another. We use it
    %  to enter Sarabun on the first Thai character and leave it on the first
    %  non-Thai one -- so Thai typed anywhere, including in \title and frame
    %  titles, just works.
    %
    %  ucharclasses must be loaded AFTER every font is set up, which is why
    %  this sits at the bottom of the file.
    % -----------------------------------------------------------------------
    \usepackage{ucharclasses}
    \setTransitionsFor{Thai}{\deckenterthai}{\deckleavethai}
  \else
    % -----------------------------------------------------------------------
    %  LuaLaTeX is NOT supported for Thai, and stops here rather than
    %  producing a broken PDF.
    %
    %  Why: automatic Thai fallback needs ucharclasses, which is built on
    %  XeTeX's \XeTeXcharclass and does not exist for LuaTeX. The obvious
    %  workaround -- make Sarabun the document-wide sans font -- makes this
    %  theme die at shipout under luaotfload with the unhelpful
    %  "error: (file ) (type 2): cannot find file ''".
    %
    %  Without either, Thai would simply vanish from the slides with no error
    %  at all, which is the exact failure this file exists to prevent. So a
    %  loud stop is the honest outcome.
    %
    %  Use XeLaTeX. It is the recommended engine for Thai anyway.
    % -----------------------------------------------------------------------
    \PackageError{lang-thai}{LuaLaTeX cannot be used for Thai in this template}{%
      Build with XeLaTeX instead:^^J%
        - line 1 of main.tex:  %% !TeX program = xelatex^^J%
        - .latexmkrc:          $pdf_mode = 5;^^J%
      pdfLaTeX also works, and falls back to the Garuda font.^^J%
      ^^J%
      Reason: automatic Thai font fallback needs ucharclasses, which is
      XeLaTeX-only.}
  \fi

  % --- Compatibility aliases ----------------------------------------------
  %  Nothing needs wrapping any more. These stay so that older slides written
  %  against the previous version of this template keep working unchanged.
  \newenvironment{thaipar}{\par}{\par}
  \newcommand{\thaiinline}[1]{#1}
  \providecommand{\textthai}[1]{#1}
\fi

% ===========================================================================
%  NOTES
% ---------------------------------------------------------------------------
%  MIXING THAI AND LATIN
%      Type them together and the font follows the script automatically:
%          The word ตัวอย่าง means ``example''.
%      Weight and shape carry across the boundary, so this stays bold
%      throughout:
%          \textbf{Results / ผลลัพธ์}
%
%  IF THAI COMES OUT BLANK
%      You are almost certainly not on XeLaTeX. Check the first line of
%      main.tex says  % !TeX program = xelatex  and that the deck's .latexmkrc
%      has  $pdf_mode = 5;
%
%      To make the blanks visible rather than silent, add this to main.tex:
%          \tracinglostchars=3
%      Every dropped character then becomes a logged error naming the font.
%
%  polyglossia IS DELIBERATELY NOT LOADED
%      It was, in an earlier version. Two reasons it is gone:
%        1. It only switches font on an explicit \textthai / \begin{thai},
%           which is exactly the wrapping this file exists to abolish.
%        2. It writes \xpg@aux into main.aux. If a later run used pdfLaTeX --
%           which has no polyglossia -- that stale .aux killed the build with
%           "Undefined control sequence \xpg@aux". Editors such as TeXstudio
%           build in place and hit this constantly. Dropping polyglossia
%           removes the failure at its source.
% ===========================================================================
